\chapter{多尺度热质耦合机理与降阶状态建模}
\label{ch:physics}

停机、探针与解释最少要读温度、液相含水和反应进度。中心温度越过某一阈值只是跟踪器的中间量。十六条控制体上的工程通量构成评分核，连续方程留在规格层。本章由连续形式收缩到实际积分所用的离散核，收缩后保留冷点、表面温度、平均含水和致死进度；压力场、蒸汽质量分数和三维流动不进积分。界面潜热由蒸汽指令与活度的工程式给出，符号与冷凝放热、蒸发吸热一致。

\section{计算域与全场状态}

食品区域 \(\Omega_f\subset\R^3\)，腔体气体区域 \(\Omega_c\)，交界面 \(\Gamma=\partial\Omega_f\)。体积分数满足
\[
\varepsilon_s+\varepsilon_l+\varepsilon_g=1,\qquad \varepsilon_s=\varepsilon_s(S)>0.
\]
\(S(\vect{r},t)\) 综合收缩、孔隙闭合与表面结壳。全场状态为
\begin{equation}
\label{eq:state}
\vect{X}(\vect{r},t)=\big(T,\,p_g,\,\varepsilon_l,\,\omega_v,\,\vect{\alpha},\,S\big).
\end{equation}
反应进度 \(\vect{\alpha}=(\alpha_p,\alpha_g,\alpha_b,\alpha_m)\) 依次是蛋白质变性、淀粉糊化、美拉德褐变与微生物致死当量。腔体状态不进入食品场，而作为边界输入 \((T_\infty,Y_\infty,\vect{v}_\infty,\dot m_{\mathrm{st}},P_{\mathrm{heat}})\)。十六条控制体并不把 \(p_g\) 和 \(\omega_v\) 列为独立未知量。写出全场，是为了标明离散项各自对应连续方程中的哪一项。

\section{能量、水分与界面相变}

忽略热局部非平衡时，固、液、气在同一代表体元内共用一个 \(T\)；这一简化在蒸烤的分钟尺度上可以接受。能量方程保留四项：
\begin{equation}
\label{eq:energy}
\begin{aligned}
(\rho c)_{\mathrm{eff}}\frac{\partial T}{\partial t}
={}&\nabla\!\cdot\!\big(k_{\mathrm{eff}}\nabla T\big)
-\sum_{i\in\{l,v,a\}}\nabla\!\cdot\!(h_i\vect{n}_i)\\
&-\dot m_{\mathrm{phase}}\,\Delta h_{\mathrm{vap}}
+q_{\mathrm{reac}}.
\end{aligned}
\end{equation}
有效容积热容与导热随 \(\varepsilon_l\) 和 \(S\) 变化：表面变干后导热下降，构成外干内生的机制。液相与蒸汽
\begin{align}
\label{eq:liquid}
\frac{\partial(\varepsilon_l\rho_l)}{\partial t}
&=-\nabla\!\cdot\!\vect{n}_l-\dot m_{\mathrm{phase}},\\
\label{eq:vapor}
\frac{\partial(\varepsilon_g\rho_g\omega_v)}{\partial t}
&=-\nabla\!\cdot\!\vect{n}_v+\dot m_{\mathrm{phase}}.
\end{align}
通量由 Darcy 渗流与 Fick 扩散合成。相变保持有限速率：
\begin{equation}
\label{eq:phase}
\begin{aligned}
\dot m_{\mathrm{phase}}
={}&K_{\mathrm{ev}}\,a_w(\varepsilon_l,T)\,
\big(p_{\mathrm{sat}}(T)-p_v\big)_+\\
&-K_{\mathrm{cd}}\,(p_v-p_{\mathrm{sat}}(T))_+.
\end{aligned}
\end{equation}
蒸制时腔体蒸汽压高，冷凝项在表面释放潜热；烤制时蒸汽压低，蒸发项吸热并留下干壳。符号仍由 \(p_v-p_{\mathrm{sat}}\) 决定，不由档位名称决定。

本次核没有压力未知量，界面潜热收缩为
\begin{equation}
\label{eq:qlat}
\begin{aligned}
q_{\mathrm{lat}}={}&18000\,u_{\mathrm{st}}(1-a_w)\\
&-9000\,(1-u_{\mathrm{st}})\max(a_w-0.35,0)\,u_{\mathrm{fan}}.
\end{aligned}
\end{equation}
蒸汽指令直接决定符号，活度只调制幅度。式子保留冷凝放热、蒸发吸热的方向，却不再保留\eqref{eq:phase}对局部蒸汽压的依赖。另有一项 Antoine 压差项：腔体侧分压是腔湿乘以腔温下的饱和压，食品侧分压是表面活度乘以表面温度下的饱和压，两侧之差再乘 \(4000\,\mathrm{W}/(\mathrm{m}^2\cdot\mathrm{bar})\)，加到潜热上。水的饱和蒸汽压
\begin{equation}
\label{eq:antoine}
\log_{10} P^\star = 5.1962 - \frac{1730.63}{T+233.426},
\end{equation}
温度取摄氏度，压力取 bar。

\section{腔体边界的标量化与一维离散}

数值核并不单独求解三维流动场，而把腔体零件收缩为两个标量——腔温 \(T_c\) 与腔湿 \(h_c\)：
\begin{align}
\dot T_c &= 0.02\,u\,(T_h-T_c)-0.01\,e\,(T_c-25),\\
\dot h_c &= 0.015\,s-0.02\,e\,h_c.
\end{align}
换热系数采用核内定律
\[
h=12+55\,u_{\mathrm{fan}}\qquad[\mathrm{W}/(\mathrm{m}^2\mathrm{K})].
\]
风机从 \(0.4\) 升到 \(1.3\) 时 \(h\) 从 \(34\) 升到 \(83.5\)。另有一张通道关联表：风速从静止升到 \(2.5\,\mathrm{m/s}\) 时，换热系数从 \(3.20\) 升到 \(12.99\)。该表全程低于核的工作点，并且没有接到十六条的边界。正文中的 Biot 数因此用核的 \(h=12+55u_{\mathrm{fan}}\) 和 \(k=0.68\,\mathrm{W}/(\mathrm{m\cdot K})\)。辐射没有单独的角系数场，加热指令上保留 \(7000u\) 作为工程份额。

食品沿厚度分成十六条控制体。下标 \(0\) 在界面一侧，下标 \(15\) 在对称面。时间步 \(\Delta t=2\,\mathrm{s}\)。导热基线 \(k=\alpha\rho c\)，其中 \(\alpha=2.0\times 10^{-7}\,\mathrm{m}^2/\mathrm{s}\)，\(\rho c=3.4\times 10^{6}\,\mathrm{J}/(\mathrm{m}^3\mathrm{K})\)，故 \(k=0.68\,\mathrm{W}/(\mathrm{m\cdot K})\)。结壳变量 \(\xi\) 在表面含水低于 \(0.45\) 时增长，增长之后表面那一条的导热乘 \(1-0.55\xi\)。含水扩散系数 \(D=2\times 10^{-9}(1+3m)\)。\(15\,\mathrm{mm}\) 的显式稳定上限约 \(2.20\,\mathrm{s}\)，选用的 \(2\,\mathrm{s}\) 低于该上限。离散方程的完整形式见附录\ref{app:discrete}。

图\ref{fig:profiles}给出一维核沿厚度的温度剖面：表面先升，冷点滞后；后期加热下降后剖面重新拉平。图\ref{fig:xt}把同一剖面沿时间展开。探针窗口落在进门后约 \(2.5\)--\(3.2\,\mathrm{min}\)，近表面留下短暂的高温条带，尚未传到冷点；横向没有第二维梯度，这是一维核的几何。图\ref{fig:paw}给出核内活度对含水。S1 的三档初含水 \(0.64/0.72/0.80\) 都落在饱和段 \(a_w\approx 0.99\)，界面潜热对含水的导数在这一段接近零。活度是含水的单调饱和函数，与公开 GAB 等温线同方向。

\begin{figure}[htbp]
\centering
\colimg{P_profiles.png}
\caption{一维核沿厚度的温度剖面。表面先升，冷点滞后；后期加热下降后剖面重新拉平}
\label{fig:profiles}
\end{figure}

\begin{figure}[htbp]
\centering
\colimg{P_field_xt.png}
\caption{厚度--时间温度场。探针窗口在近表面留下短暂高温条带，尚未传到冷点}
\label{fig:xt}
\end{figure}

\begin{figure}[htbp]
\centering
\colimg{P_aw.png}
\caption{核内活度。竖线是三档初含水，三者都在饱和段}
\label{fig:paw}
\end{figure}

\section{进度、表面品质与结构演化}

速率形式为
\begin{equation}
\label{eq:kinetics}
\frac{\partial\alpha_i}{\partial t}
=A_i\exp\!\Big(-\frac{E_{a,i}}{RT}\Big)
(1-\alpha_i)^{n_i}\,g_i(\varepsilon_l,S).
\end{equation}
\(g_p\) 在自由水不足时下降，\(g_b\) 在低水分、高温时上升，使表面褐变和内部变性可以分离。热力进度不另造机理：
\begin{equation}
\label{eq:F}
\frac{\mathrm{d}F}{\mathrm{d}t}=10^{(T_{\mathrm{cold}}-T_{\mathrm{ref}})/z},\qquad
\alpha_m=1-\exp(-F/F_{\mathrm{req}}).
\end{equation}
冷点取食品域上温度场的下确界 \(T_{\mathrm{cold}}(t)=\min_{\vect{r}\in\Omega_f}T(\vect{r},t)\)。核取十六条控制体的最低温度。\(\alpha_m=1-\mathrm{e}^{-F/2}=0.85\) 给出 \(F=-2\ln 0.15=3.794\,\mathrm{min}\)。冷点若保持恒温 \(T\)，每分钟累积 \(10^{(T-70)/10}\) 分钟的等效时间。表\ref{tab:constT}只使用这条映射，不使用任何一条轨迹。

设计线是策略比较的零点，不是认证阈值。面板若把停机表述为“中心到达 \(70^\circ\mathrm{C}\)”，给出的只是跟踪器的中间量：同一条 \(70^\circ\mathrm{C}\) 恒温，累积 \(F\) 的速度是每分钟 \(1\)，走到 \(0.85\) 需要 \(3.794\,\mathrm{min}\)；若冷点只到 \(60^\circ\mathrm{C}\)，每分钟只累积 \(0.1\)，走完同一条线要将近 \(38\,\mathrm{min}\)。微生物初始数量不进入状态。

\begin{table}[htbp]
\centering
\caption{恒温冷点到达 \(\alpha_m=0.85\) 的时间}
\label{tab:constT}
{\small
\begin{tabular}{@{}rrrr@{}}
\toprule
\(T/^\circ\mathrm{C}\) & 每分钟的 \(F\) & 所需时间 / min & 所需时间 / s \\
\midrule
55 & 0.0316 & 119.98 & 7199 \\
60 & 0.100 & 37.94 & 2277 \\
65 & 0.316 & 12.00 & 720 \\
70 & 1 & 3.794 & 228 \\
75 & 3.162 & 1.200 & 72 \\
80 & 10 & 0.379 & 23 \\
\bottomrule
\end{tabular}}
\end{table}

结构演化
\begin{equation}
\label{eq:structure}
\frac{\partial S}{\partial t}
=\kappa_s(1-S)\,\langle\alpha_p-\alpha_s^\star\rangle_+
+\kappa_c\,\langle M_{c}- \varepsilon_l\rangle_+.
\end{equation}
核里结壳只在表面含水低于 \(0.45\) 时以 \(1.5\times 10^{-4}\,\mathrm{s}^{-1}\) 增长。S1 的初始含水不低于 \(0.64\)，表面失水率很低，结壳常停在 \(0\)。规格中的路径 C 因此在多数样本上没有闭合。褐变 \(\alpha_b\) 要求表面温度高于 \(110^\circ\mathrm{C}\) 且表面含水低于 \(0.55\)；家用组合档把表面湿度维持得较高，褐变通道几乎不动作。

三个原型共用同一套导热和进度公式，差别只在活度尺度：肉糜 \((0.10,0.70)\)，土豆 \((0.055,0.80)\)，面团 \((0.08,0.75)\)。交叉表里三个食材的中位误差几乎相同，与“原型只改了一个弱的活度尺度”一致；该重合只说明活度尺度的扰动不足以推动中位误差，不能当作蛋白与淀粉反应差异已经过检验的证据。面团膨胀不在状态里。

Kossovich 数
\[
\mathrm{Ko}
=\frac{\Delta h_{\mathrm{vap}}\,\varepsilon_{l0}\,\rho_l}{(\rho c)_{\mathrm{eff}}\,\Delta T}
\]
停在两个未定因子上。核文件里写着潜热常数 \(\Delta h_{\mathrm{vap}}=2.3\times 10^{6}\,\mathrm{J/kg}\)。步进不读取这个常数：界面潜热走的是 \(18000\,u_{\mathrm{st}}\) 与 \(9000\,u_v\) 两条工程通量。体积热容 \((\rho c)_{\mathrm{eff}}=3.4\times 10^{6}\,\mathrm{J}/(\mathrm{m}^3\mathrm{K})\) 进入了温升。两者之比是 \(0.6765\,\mathrm{K\cdot m}^3/\mathrm{kg}\)，于是 \(\mathrm{Ko}=0.6765\,\varepsilon_{l0}\rho_l/\Delta T\)。核里的含水是夹在 \([0.02,0.95]\) 的分数，不是液相体积分数乘密度。因此 Ko 留在这两个因子上，不化成一个纯数。

\section{两节点模型与内部热阻}

探针分析使用表面节点 \(1\) 与内部节点 \(2\)：
\begin{align}
\label{eq:twonode}
C_1\dot T_1
&=\frac{T_\infty-T_1}{R_{\mathrm{ext}}(v)}+\frac{T_2-T_1}{R_{\mathrm{int}}}
+q_{\mathrm{lat}}(\dot m_\Gamma),\\
C_2\dot T_2
&=\frac{T_1-T_2}{R_{\mathrm{int}}},\\
\dot m
&=-k_m\big(a_w(m)\,p_{\mathrm{sat}}(T_1)-p_\infty\big).
\end{align}
\(R_{\mathrm{ext}}\) 由风速与是否冷凝切换，\(R_{\mathrm{int}}\) 随结壳增大。无相变且 \(\dot m=0\) 时，平衡迫使 \(T_1=T_2=T_\infty\)，三个参数从代数方程中消失。把含水冻结，相变通量为零，对 \((T_1,T_2)\) 线性化，
\[
\vect{A}=\begin{pmatrix}-a-b & b\\ c & -c\end{pmatrix},\qquad
\vect{C}=\begin{pmatrix}1 & 0\end{pmatrix},
\]
其中 \(a=1/(C_1 R_{\mathrm{ext}})\)，\(b=1/(C_1 R_{\mathrm{int}})\)，\(c=1/(C_2 R_{\mathrm{int}})\)。可观测矩阵行列式等于 \(b\neq 0\)，两个温度状态从表面看是可观测的。平衡点上对 \(R_{\mathrm{int}}\) 的灵敏度为零。状态可观测只说明温度节点可以从表面重建；内部热阻的辨识还需要离开平衡的输入，让灵敏度列不再是零。

说明常数 \(C_1=C_2=2000\,\mathrm{J/K}\)，\(R_{\mathrm{ext}}=0.05\,\mathrm{K/W}\)，\(R_{\mathrm{int}}=0.15\,\mathrm{K/W}\) 时，特征根对应时间常数 \(69.72\,\mathrm{s}\) 与 \(430.3\,\mathrm{s}\)。\(40\,\mathrm{s}\) 的探针窗短于这两个时间常数：表面记录在代数上满秩，数值上仍然弱，窗口结束时冷点几乎还停在初值，信息主要落在表面节点的瞬态上。

评分循环始终积分十六条控制体。压力场、蒸汽质量分数和有限速率相变源项写在连续形式里，用来标明离散项各自对应哪一项；核里实际推进的是温度、含水、结壳代理和进度。两节点工作方程只服务于探针分析：它给出平衡点上内部热阻灵敏度为零、两段脉冲提供第二个灵敏度方向这两条代数事实。

\section{过程特征尺度与二维冷点偏移}

半厚取 \(L=30\,\mathrm{mm}\)，Fourier 数 \(\mathrm{Fo}=\alpha t/L^2\)，Biot 数 \(\mathrm{Bi}=hL/k\)。表\ref{tab:fobi}给出设计尺度。探针窗的 Fourier 数为 \(8.9\times 10^{-3}\)，远小于 \(1\)，冷点来不及跟上表面。探针段把 Biot 数从 \(1.50\) 升到 \(3.68\)，表面梯度更陡，内部更跟不上。这是两段探针要把外部热阻和内部热阻分开的尺度理由：第一段改相变符号，第二段改 \(R_{\mathrm{ext}}\)，内部节点在 \(40\,\mathrm{s}\) 里几乎不动。图\ref{fig:fobi}把 \(\mathrm{Fo}(t)\) 与 \(\mathrm{Bi}(u_{\mathrm{fan}})\) 画在核常数下。

\begin{table}[htbp]
\centering
\caption{核常数下的 Fourier 数与 Biot 数。\(L=30\,\mathrm{mm}\)}
\label{tab:fobi}
\setlength{\tabcolsep}{3.2pt}
{\small
\begin{tabular}{lrrr}
\toprule
条件 & \(t\) 或 \(u_{\mathrm{fan}}\) & \(\mathrm{Fo}\) 或 \(h\) & \(\mathrm{Bi}\) \\
\midrule
探针窗 & \(40\,\mathrm{s}\) & \(8.9\times 10^{-3}\) & --- \\
开环 \(30\,\mathrm{mm}\) 菜单 & \(1170\,\mathrm{s}\) & \(0.26\) & --- \\
家用时域终点 & \(3300\,\mathrm{s}\) & \(0.73\) & --- \\
风机 \(0.4\) & \(h=34\) & --- & \(1.50\) \\
探针段风机 \(1.3\) & \(h=83.5\) & --- & \(3.68\) \\
\bottomrule
\end{tabular}}
\end{table}

\begin{figure*}[t]
\centering
\spanimg{X_fo_bi.png}
\caption{核常数下的 Fourier 数随时间、Biot 数随风机。探针窗落在 \(\mathrm{Fo}\ll 1\)、\(\mathrm{Bi}>1\) 的区域}
\label{fig:fobi}
\end{figure*}

一维核把冷点取成十六条中的最低温度，并默认冷点落在对称面。门缝漏汽、搁架遮挡、一侧靠壁、脂肪夹层偏下，都会把最低温度推离几何中心。一块独立的二维显式截面沿厚度 \(24\) 格、沿横向 \(16\) 格，时间 \(20\,\mathrm{min}\)，边界固定。脂肪夹层关闭时，\(20\,\mathrm{min}\) 终点冷点 \(36.18^\circ\mathrm{C}\)，表面 \(128.62^\circ\mathrm{C}\)，最低温度落在绝热侧格子 \(23\)，几何中心在格子 \(12\)，偏移 \(13.75\,\mathrm{mm}\)。脂肪夹层打开时，终点冷点 \(34.11^\circ\mathrm{C}\)，表面 \(126.0^\circ\mathrm{C}\)，最低温度仍在绝热侧，偏移同样是 \(13.75\,\mathrm{mm}\)。图\ref{fig:t2grid}给出 \(20\,\mathrm{min}\) 的温度场。截面证明冷点可以不在中心：最低温度被绝热边界固定在最内侧。评分循环仍用一维核的十六条，是因为交叉表比较的是策略，不是搁架几何。把二维偏移写进一维冷点定义，会把“最低温度所在格”改成“横向再取一次下确界”，主表没有这一步。

\begin{figure}[htbp]
\centering
\colimg{X_t2_grid.png}
\caption{\(24\times 16\) 显式截面在 \(20\,\mathrm{min}\) 的温度场。最低温度在绝热侧，偏离中心 \(13.75\,\mathrm{mm}\)}
\label{fig:t2grid}
\end{figure}

交叉表、网格加密和典型轨迹都读取同一组工程常数：扩散率 \(\alpha=2.0\times 10^{-7}\,\mathrm{m}^2/\mathrm{s}\)，容积热容 \(\rho c=3.4\times 10^{6}\,\mathrm{J}/(\mathrm{m}^3\mathrm{K})\)，基线导热 \(k=0.68\,\mathrm{W}/(\mathrm{m\cdot K})\)，时间步 \(2\,\mathrm{s}\)，时域 \(55\,\mathrm{min}\)，控制体数 \(16\)；换热 \(h=12+55u_{\mathrm{fan}}\)，界面潜热按\eqref{eq:qlat}，辐射份额 \(7000u\)，Antoine 压差再乘 \(4000\,\mathrm{W}/(\mathrm{m}^2\cdot\mathrm{bar})\)；脂肪夹层把指定控制体的导热乘 \(0.45\)。这些数字一旦改写，已经写出的日志不再对应。
